% Self-contained section for integration into the parent research article.
% Uses standard theorem/lemma/corollary/proposition environments and amsmath.
\section{The optimal exterior differential barrier}
\label{opt:sec:barrier}

The incoming report \emph{Lerch Global Phase} proves an exterior
one-crossing inequality with the fixed coefficient $e^{-n}$ and explicitly
asks for the optimum over all positive coefficients in its further-research
question~2. We solve that optimization problem. The sharp coefficient is
controlled by a cubic algebraic number; its large-index limit involves the
golden ratio. A distinction between a strict pointwise inequality and a
strict inequality for positive geometric sums is essential at the optimum.

For integers $n\geq2$ define
\begin{equation}\label{opt:eq:family}
 g_n(x)=\frac{(\log x)^{n-1}(n-\log x)}{x^2},\qquad
 U_n(a,\rho)=\sum_{m=0}^{\infty}\rho^m g_n(a+m).
\end{equation}
Here $a>0$ and $0\leq\rho\leq1$. The series and each fixed number
of its $a$ derivatives converge absolutely and locally uniformly,
uniformly in $\rho$, since the $m$th summand after $j$ derivatives is
$O(m^{-2-j}(1+\log^n m))$. For $\rho<1$, $U_n$ is the first
$a$ derivative of the order-$n$ Lerch coefficient
$\sum_{m\geq0}\rho^m\log^n(a+m)/(a+m)$; at $\rho=1$ it is
$\gamma_n'(a)$ in the usual Stieltjes Laurent convention.

Call a pair $(B,\lambda)$ \emph{uniformly admissible} if $B>0$,
$\lambda>0$, and
\begin{equation}\label{opt:eq:admissible}
 \partial_aU_n(a,\rho)+\lambda U_n(a,\rho)<0
 \qquad(a\geq B,\ 0<\rho\leq1).
\end{equation}
The uniformity down to arbitrarily small positive $\rho$ makes this a
precise optimization problem. A smaller cutoff at one fixed deformation
parameter is a different problem.

Put $\varphi=(1+\sqrt5)/2$. Let $r_n$ be the unique positive root of
\begin{equation}\label{opt:eq:cubic}
 r^3+nr^2-nr-n=0,
\end{equation}
and set
\begin{equation}\label{opt:eq:lambda}
 \tau_n=(2-r_n)e^{-r_n},\qquad \lambda_n=e^{-n}\tau_n,
 \qquad \Delta_n=\sqrt{n^2+8n},\qquad
 \alpha_n=\frac{3n-\Delta_n}{4}.
\end{equation}
For $0<\tau\leq1$, write
\begin{equation}\label{opt:eq:R}
 R_{n,\tau}(y)=2y^2-3ny+n(n-1)+\tau e^{y-n}y(n-y).
\end{equation}

\begin{theorem}[Optimal coefficient and optimal cutoff]
\label{opt:thm:optimal}
For every integer $n\geq2$ the following statements hold.
\begin{enumerate}
\item The cubic root satisfies $1<r_n<\varphi$. There is a unique
 $\eta_n\in(\alpha_n,(n-1)/2)$ such that
 $R_{n,\tau_n}(\eta_n)=0$. Put $B_n=e^{\eta_n}$.
\item The pointwise inequality is
 \begin{equation}\label{opt:eq:pointwise}
  g_n'(x)+\lambda_n g_n(x)\leq0\qquad(x\geq B_n).
 \end{equation}
 Equality occurs at exactly two points,
 \begin{equation}\label{opt:eq:equality-points}
  x=B_n\quad\hbox{and}\quad x=e^{n+r_n}.
 \end{equation}
\item The pair $(B_n,\lambda_n)$ is uniformly admissible. Every
 uniformly admissible pair $(B,\lambda)$ has $B\geq B_n$, and if
 $B=B_n$ then $\lambda=\lambda_n$.
\item Consequently $e^{\lambda_n a}U_n(a,\rho)$ is strictly decreasing
 on $[B_n,\infty)$ for every $0<\rho\leq1$. There is at most one
 zero there; any such zero is simple and crosses from positive to
 negative.
\end{enumerate}
\end{theorem}

We give a complete proof, including the boundary qualification.

\begin{lemma}[The negative-tail constraint]\label{opt:lem:tail}
For $r>0$ put
\[
 T_n(r)=2-\frac{n-1}{n+r}-\frac1r.
\]
The function $e^{-r}T_n(r)$ has a unique maximum, at $r=r_n$, and
\begin{equation}\label{opt:eq:Tmaximum}
 \max_{r>0}e^{-r}T_n(r)=(2-r_n)e^{-r_n}=\tau_n.
\end{equation}
Thus $\lambda_n$ is the smallest coefficient that can make
$g_n'+\lambda g_n\leq0$ on the whole interval $(e^n,\infty)$.
\end{lemma}
\begin{proof}
For $x=e^{n+r}>e^n$ one has $g_n(x)<0$ and
\[
 -\frac{g_n'(x)}{g_n(x)}=e^{-n-r}T_n(r).
\]
Dividing the desired inequality by the negative quantity $g_n(x)$
shows that it is equivalent to
$\lambda\geq e^{-n-r}T_n(r)$ for every $r>0$.

Now
\[
 T_n'(r)=\frac{n-1}{(n+r)^2}+\frac1{r^2}>0,
 \qquad
 T_n''(r)=-\frac{2(n-1)}{(n+r)^3}-\frac2{r^3}<0.
\]
Therefore $T_n'-T_n$ is strictly decreasing, tends to $+\infty$
at zero, and tends to $-2$ at infinity. It has precisely one zero,
which is the unique maximum of $e^{-r}T_n(r)$. The factorization
\begin{equation}\label{opt:eq:factor}
 (T_n-T_n')r^2(n+r)^2
 =(n+2r)(r^3+nr^2-nr-n)
\end{equation}
identifies that point with the positive root of \eqref{opt:eq:cubic}.
The cubic has exactly one positive root by its coefficient sign
variation, or by the preceding uniqueness argument. Its values at $1$
and $\varphi$ are $1-n<0$ and $\varphi^3>0$, respectively, giving
$1<r_n<\varphi$.

Rearranging the cubic gives
$n=r_n^3/(1+r_n-r_n^2)$. Substitution in $T_n(r_n)$ simplifies it
to $2-r_n$, proving \eqref{opt:eq:Tmaximum}. The maximum is positive;
the tail inequality is strict except at its unique maximizing point.
\end{proof}

\begin{lemma}[The positive-side cutoff]\label{opt:lem:cutoff}
For every $n\geq2$ and every $0<\tau\leq1$, the function
$R_{n,\tau}$ has exactly one root $\eta_n(\tau)$ in $(0,n)$.
It satisfies
\[
 \alpha_n<\eta_n(\tau)<\frac{n-1}{2},
 \qquad
 R_{n,\tau}(y)>0\ (0<y<\eta_n(\tau)),\qquad
 R_{n,\tau}(y)<0\ (\eta_n(\tau)<y\leq n).
\]
The root is simple, and $\eta_n(\tau)$ is strictly increasing in $\tau$.
\end{lemma}
\begin{proof}
For $y=n-t$, $0\leq t\leq(n+1)/2$, use
$\tau e^{-t}\leq(1+t)^{-1}$ and $(n-t)t\geq0$. This gives
\begin{align*}
 (1+t)R_{n,\tau}(n-t)
 &\leq2t^3+(1-n)t^2-nt-n\\
 &\leq2t^2-nt-n<0.
\end{align*}
For the second inequality use $2t\leq n+1$. The final quadratic is
convex, with negative values $-n$ and $(1-n)/2$ at the two endpoints,
so it is negative throughout. This proves negativity for
$(n-1)/2\leq y\leq n$.

On $0\leq y\leq(n-1)/2$,
\[
 R_{n,\tau}''(y)
 =4+\tau e^{y-n}\{-y^2+(n-4)y+2n-2\}>4.
\]
The bracket is a concave quadratic whose endpoint values are
$2(n-1)$ and $(n^2-1)/4$, both positive. Since $R_{n,\tau}(0)>0$
and $R_{n,\tau}((n-1)/2)<0$, strict convexity gives one and only one
crossing on this interval; a second crossing would force the function
to be nonnegative at its right endpoint. The derivative at the crossing
is strictly negative. The polynomial part vanishes at $\alpha_n$, so
$R_{n,\tau}(\alpha_n)>0$, proving the lower bound. Finally
\[
 \frac{\partial\eta_n(\tau)}{\partial\tau}
 =-\frac{e^{\eta_n(\tau)-n}\eta_n(\tau)(n-\eta_n(\tau))}
 {\partial_yR_{n,\tau}(\eta_n(\tau))}>0.
\]
\end{proof}

\begin{proof}[Proof of Theorem~\ref{opt:thm:optimal}]
For $x=e^y>1$, differentiation gives
\begin{equation}\label{opt:eq:pointwise-factor}
 g_n'(x)+\tau e^{-n}g_n(x)
 =x^{-3}y^{n-2}R_{n,\tau}(y).
\end{equation}
The prefactor is positive. Because $1<r_n<\varphi<2$,
$0<\tau_n<1$. Lemma~\ref{opt:lem:cutoff} treats the region below
$e^n$, Lemma~\ref{opt:lem:tail} treats the region above it, and
at $x=e^n$ one has $g_n'<0$. This proves the nonpositive inequality
and the exact equality set.

For $a\geq B_n$, every summand in
\[
 (\partial_a+\lambda_n)U_n(a,\rho)
 =\sum_{m=0}^{\infty}\rho^m
       [g_n'(a+m)+\lambda_n g_n(a+m)]
\]
is nonpositive. Only two possible arguments can give equality.
When $\rho>0$, at least one of the three summands $m=0,1,2$ is
strictly negative, so the sum is strictly negative. Absolute convergence
justifies the differentiation and summation, also at $\rho=1$.

To prove optimality, suppose first that $\lambda<\lambda_n$.
At $x_*=e^{n+r_n}$, where $g_n(x_*)<0$, one has
$g_n'(x_*)+\lambda g_n(x_*)>0$. Since
$(\partial_a+\lambda)U_n(a,\rho)$ tends to this pointwise expression
as $\rho\downarrow0$, every cutoff $B\leq x_*$ fails for small
positive $\rho$. In particular, such a coefficient cannot beat $B_n$.
If $\lambda\geq\lambda_n$ and $B<B_n$, choose
$\max(B,1)<a<B_n$. Then $g_n(a)>0$ and
\[
 g_n'(a)+\lambda g_n(a)
 \geq g_n'(a)+\lambda_n g_n(a)>0,
\]
again contradicting admissibility at sufficiently small positive $\rho$.
This proves $B\geq B_n$. At $B=B_n$, any $\lambda>\lambda_n$
makes the pointwise expression strictly positive at $a=B_n$,
so only $\lambda_n$ is possible. The integrating factor and the sign
of the derivative at any zero prove the last assertion.
\end{proof}

\paragraph{Why the pointwise strict optimum is not attained.}
If the requirement is instead
$g_n'(x)+\lambda g_n(x)<0$ for every $x>B$, the infimum of possible
cutoffs is still $B_n$, but no coefficient attains it. At
$\lambda_n$ the outer point $e^{n+r_n}>B_n$ has equality.
Coefficients $\lambda>\lambda_n$, sufficiently close to $\lambda_n$,
have strict tail inequality and cutoff $e^{\eta_n(e^n\lambda)}$
decreasing to $B_n$ as $\lambda\downarrow\lambda_n$.
The strict summed theorem therefore must not be copied verbatim into
a statement about every positive measure or about the elementary
endpoint $\rho=0$.

\paragraph{Positive measures.}
Let $\nu$ be a nonzero positive measure on $[0,\infty)$ satisfying
\[
 \int_{[0,\infty)}\frac{1+\log^n(2+t)}{(1+t)^2}\,d\nu(t)<\infty.
\]
This condition supplies locally uniform domination of the integral
and its derivative. The function
$H_n(a)=\int g_n(a+t)\,d\nu(t)$ satisfies
$H_n'+\lambda_nH_n\leq0$ for $a\geq B_n$. It is strictly negative
unless the measure is supported on the at most two points for which
$a+t$ belongs to \eqref{opt:eq:equality-points}. In particular,
non-atomic measures and measures with at least three support points
give the strict inequality. This is the exact measure-theoretic
qualification forced by the optimal tangency.

\subsection{The golden-ratio limit and convergent coefficient formulas}
\label{opt:sec:asymptotic}

\begin{theorem}[Monotone optimizer and its large-index expansion]
\label{opt:thm:asymptotic}
The sequence $r_n$ strictly increases to $\varphi$, whereas
$\tau_n=e^n\lambda_n$ strictly decreases to
$\tau_\infty=(2-\varphi)e^{-\varphi}$. As $n\to\infty$,
\begin{align}
 r_n={}&\varphi-
 \left(1+\frac{2\sqrt5}{5}\right)\frac1n
 +\left(\frac52+\frac{57\sqrt5}{50}\right)\frac1{n^2}
 +O(n^{-3}),\label{opt:eq:r-asymptotic}\\
 \lambda_n={}&e^{-n}\tau_\infty
 \left[1+\frac{7+3\sqrt5}{2n}
 -\frac{35+16\sqrt5}{10n^2}+O(n^{-3})\right].
 \label{opt:eq:lambda-asymptotic}
\end{align}
Writing $\epsilon=1/n$, the complete expansion of $r_n$ is a
convergent power series for $|\epsilon|<5/27$. Its coefficients are
\begin{equation}\label{opt:eq:r-Lagrange}
 r(\epsilon)=\varphi+
 \sum_{j\geq1}\frac{(-\epsilon)^j}{j}
 [z^{j-1}]
 \left(\frac{(\varphi+z)^3}{\sqrt5+z}\right)^j.
\end{equation}
In particular this is a convergent formula for every integer $n\geq6$.
\end{theorem}
\begin{proof}
The identity
$n=r^3/(1+r-r^2)$ has strictly positive derivative on $1<r<\varphi$:
\[
 \frac{d}{dr}\frac{r^3}{1+r-r^2}
 =\frac{r^2(3+2r-r^2)}{(1+r-r^2)^2}>0.
\]
It tends to infinity at $\varphi$, proving the claims about $r_n$.
Since $d[(2-r)e^{-r}]/dr=(r-3)e^{-r}<0$ in this interval,
the claims about $\tau_n$ follow.

The implicit equation
$\epsilon r^3+r^2-r-1=0$ has derivative $\sqrt5$ in $r$ at
$(\epsilon,r)=(0,\varphi)$, so the analytic implicit-function theorem
applies. Substitution of a quadratic series gives
\eqref{opt:eq:r-asymptotic}, and substitution into $(2-r)e^{-r}$ gives
\eqref{opt:eq:lambda-asymptotic}. With $r=\varphi+z$ the equation becomes
\[
 z=-\epsilon\frac{(\varphi+z)^3}{\sqrt5+z}.
\]
Lagrange inversion gives \eqref{opt:eq:r-Lagrange}.

For completeness the convergence radius can be identified exactly.
The discriminant in $r$ is
\[
 \operatorname{disc}_r(\epsilon r^3+r^2-r-1)
 =(1-\epsilon)(27\epsilon+5).
\]
The branch selected at zero is holomorphic throughout
$|\epsilon|<5/27$: away from zero the discriminant is nonzero there,
and the selected branch has a removable, already analytic continuation
at zero. Along the real interval $-5/27<\epsilon\leq0$ this branch
increases from $\varphi$ to $3$. At $(\epsilon,r)=(-5/27,3)$,
the first $r$ derivative vanishes while
$\partial_\epsilon=27$ and $\partial_r^2=-4/3$ are nonzero.
It is therefore a genuine square-root branch point. Hence the Taylor
radius is exactly $5/27$.
\end{proof}

\subsection{The exponentially small displacement of the cutoff}
The optimizer has an ordinary convergent expansion in $1/n$. The
cutoff also contains an exponentially small displacement in logarithmic
coordinates, which should be kept separate from that algebraic expansion.
Put
\[
 E_n=e^{\alpha_n-n},\qquad A_n=\alpha_n(n-\alpha_n),
 \qquad C_n=n-2\alpha_n.
\]

\begin{theorem}[Cutoff expansion and an all-orders formula]
\label{opt:thm:cutoff-expansion}
As $n\to\infty$,
\begin{align}
 \eta_n={}&\alpha_n+
 \frac{\tau_nA_n}{\Delta_n}E_n\notag\\
 &+\tau_n^2\left[
 \frac{A_n(A_n+C_n)}{\Delta_n^2}
 +\frac{2A_n^2}{\Delta_n^3}\right]E_n^2
 +O(n^3E_n^3).\label{opt:eq:eta-two-sectors}
\end{align}
Every further power of $E_n$ has an explicit coefficient. More precisely,
for all sufficiently large $n$ the solution $\delta=\eta_n-\alpha_n$
is obtained, at $E=E_n$, from the locally convergent series
\begin{equation}\label{opt:eq:eta-Lagrange}
 \delta(E)=\sum_{j\geq1}\frac{E^j}{j}[z^{j-1}]
 \left[
 \frac{\tau_ne^z(\alpha_n+z)(n-\alpha_n-z)}{\Delta_n-2z}
 \right]^j.
\end{equation}
For each fixed truncation order $J$, its error at $E_n$ is
$O_J(n^{J+1}E_n^{J+1})$.
In the original $a$ coordinate this implies
\begin{equation}\label{opt:eq:B-asymptotic}
 B_n=e^{\alpha_n}+
 \frac{\tau_\infty}{4e^2}
 \left(n+\frac{7+3\sqrt5}{2}\right)+O(n^{-1}).
\end{equation}
The leading additive displacement above the first stationary point
of $g_n$ has coefficient
\[
 \frac{\tau_\infty}{4e^2}=0.0025625510827\ldots.
\]
\end{theorem}
\begin{proof}
Substitute $\eta_n=\alpha_n+\delta$ in its defining equation and use
the exact quadratic Taylor expansion of the polynomial part:
\[
 -\Delta_n\delta+2\delta^2+
 \tau_n E_ne^\delta(\alpha_n+\delta)(n-\alpha_n-\delta)=0.
\]
Equivalently, $\delta=E_nf_n(\delta)$ with
\[
 f_n(z)=
 \frac{\tau_ne^z(\alpha_n+z)(n-\alpha_n-z)}{\Delta_n-2z}.
\]
This proves \eqref{opt:eq:eta-Lagrange} by Lagrange inversion. Its
first two coefficients are $f_n(0)$ and $f_n(0)f_n'(0)$, giving
\eqref{opt:eq:eta-two-sectors}.

The remainder statement can be justified uniformly in $n$. On a fixed
small complex disk in $z$, $f_n$ and its derivative are $O(n)$, while
$\Delta_n\asymp n$. For $|E|<c/n$ with a sufficiently small absolute
$c>0$, the map $z\mapsto Ef_n(z)$ is a contraction of a smaller
fixed disk. The analytic solution is uniformly bounded there. Cauchy
estimates therefore bound its $j$th coefficient by $O(C^jn^j)$.
Since $E_n\asymp e^{-n/2}$, every fixed-order tail has the asserted
bound for sufficiently large $n$.

Finally exponentiate the first sector. Its omitted terms become
$O(n^2e^{-n/2})$ after multiplication by $e^{\alpha_n}$. Direct
expansion gives
\[
 \frac{A_n}{\Delta_n}e^{2\alpha_n-n}
 =e^{-2}\left(\frac n4-\frac1n+O(n^{-2})\right).
\]
Combining this with \eqref{opt:eq:lambda-asymptotic} yields
\eqref{opt:eq:B-asymptotic}. No relative approximation to
$\eta_n-\alpha_n$ is being inferred from an uncontrolled algebraic
remainder in $\alpha_n$.
\end{proof}

\paragraph{Limits of the conclusion.}
The optimum is over constant first-order operators, required uniformly
for all positive geometric weights down to $\rho=0$. It does not
optimize variable-coefficient operators, inequalities only at the zeros,
or cutoffs at one fixed $\rho$. Those extensions remain separate
research questions. In particular, improving an exterior barrier does
not on its own classify all interior zeros or prove a global
monotonicity theorem for individual zero branches.
